\renewcommand{\thetheorem}{6.\arabic{theorem}}
\setcounter{theorem}{0}
\setcounter{chapter}{5}
% ---------------------------------------------------------------------------
% Provenance of the prose in this file.
%   % Extractor: <model>   the block below was transcribed from Prof. Biran's
%                          handwritten notes by that model
%   % Creator:   <model>   the block below was written by that model and has no
%                          counterpart in the notes
% A tag holds until the next tag.
%
% This chapter was written before the convention existed, so the prose carries
% no tags except on the blocks added by later passes. Untagged prose is the
% original transcription; the AI-Notes at the end record which figures and
% expansions are the transcript's rather than the notes'.
% ---------------------------------------------------------------------------

\chapter{Systems of Linear Equations}

\subsection{Fields (Körper)}

Before we solve systems of linear equations, we must fix the algebraic structure over which these equations are built. Throughout this course, we will work over a base \newterm{field} (in German: \qt{Körper}), usually denoted by $K$.

A field is a set equipped with two operations (addition $+$ and multiplication $\cdot$) satisfying the ten axioms \textbf{(K1)}--\textbf{(K10)} of \cref{def:field}: commutativity, associativity and distributivity, together with the existence of the identities $0$ and $1$ and of additive and multiplicative inverses. We briefly recall the examples we will use most often.

\begin{example}[Common Fields]
\begin{itemize}
    \item $\mathbb{Q} = \left\{a \mid a = \frac{m}{n}, \; m, n \in \mathbb{Z}, \, n \ne 0\right\}$ (the rational numbers)
    \item $\mathbb{R}$ (the real numbers)
    \item $\mathbb{C} = \{z \mid z = a + i \cdot b, \; a, b \in \mathbb{R}\}$, where $i \cdot i = -1$ (the complex numbers)
\end{itemize}
\end{example}

\begin{example}[Finite Fields]
Fields do not have to be infinite. We will frequently use finite fields, where arithmetic is performed modulo a prime number.
\begin{enumerate}
    \item $\mathbb{F}_2 = \{0, 1\}$. Operations are done modulo $2$:
    \[
    \begin{array}{c|cc}
    + & 0 & 1 \\
    \hline
    0 & 0 & 1 \\
    1 & 1 & 0
    \end{array}
    \qquad\qquad
    \begin{array}{c|cc}
    \cdot & 0 & 1 \\
    \hline
    0 & 0 & 0 \\
    1 & 0 & 1
    \end{array}
    \]
    \item $\mathbb{F}_3 = \{0, 1, 2\}$. Operations are done modulo $3$:
    \[
    \begin{array}{c|ccc}
    + & 0 & 1 & 2 \\
    \hline
    0 & 0 & 1 & 2 \\
    1 & 1 & 2 & 0 \\
    2 & 2 & 0 & 1
    \end{array}
    \qquad\qquad
    \begin{array}{c|ccc}
    \cdot & 0 & 1 & 2 \\
    \hline
    0 & 0 & 0 & 0 \\
    1 & 0 & 1 & 2 \\
    2 & 0 & 2 & 1
    \end{array}
    \]
\end{enumerate}
\end{example}

\subsection{A Motivating Example}

Let us solve a small system of linear equations over $\mathbb{R}$ by elimination:
\begin{equation*}
\begin{cases}
R_1: \quad x + 3y + 4z = -1 \\
R_2: \quad 2x - \phantom{3}y + \phantom{4}z = \phantom{-}5
\end{cases}
\end{equation*}

We operate on the equations so as to eliminate variables one at a time. First, we remove $x$ from the second equation by $R_2 - 2 \cdot R_1 \to R_2$:
\begin{equation*}
\begin{cases}
x + 3y + 4z = -1 \\
\phantom{x} \; -7y - 7z = \phantom{-}7
\end{cases}
\end{equation*}
Next, we normalise the second equation by $-\frac{1}{7} \cdot R_2 \to R_2$:
\begin{equation*}
\begin{cases}
x + 3y + 4z = -1 \\
\phantom{x + 3}y + \phantom{4}z = -1
\end{cases}
\end{equation*}
Finally, instead of substituting backwards, we eliminate $y$ \emph{upwards} as well, using $R_1 - 3 \cdot R_2 \to R_1$:
\begin{equation*}
\begin{cases}
x \phantom{{}+ 3y} + \phantom{4}z = \phantom{-}2 \\
\phantom{x + 3}y + \phantom{4}z = -1
\end{cases}
\end{equation*}

The system is now solved as far as it can be. There are three unknowns but only two equations, so one unknown remains free: let $z := c \in \mathbb{R}$ be an arbitrary real number. Then
\[
y = -1 - c, \qquad x = 2 - c,
\]
and the set of solutions is
\begin{equation} \label{eq:motivating_solution_set}
\{(x, y, z) \mid x = 2 - c, \; y = -1 - c, \; z = c, \; c \in \mathbb{R}\}.
\end{equation}
This is not a single point but an entire one-parameter family, geometrically, the line in which the two planes $x + 3y + 4z = -1$ and $2x - y + z = 5$ intersect, as sketched in \cref{fig:two_planes}.

\begin{figure}[ht!]
  \centering
  \begin{tikzpicture}[scale=0.92, font=\small]
    % Two parallelograms sheared in opposite directions, both containing the
    % horizontal segment y = 0: that segment is their line of intersection.
    % The fills are translucent so that the overlap region stays visible.
    \fill[ThemeSecondary!50, fill opacity=0.42, draw=ThemeSecondary!75, thin]
      (-5.0,1.5) -- (1.8,1.5) -- (5.0,-1.5) -- (-1.8,-1.5) -- cycle;
    \fill[ThemeGreen!50, fill opacity=0.42, draw=ThemeGreen!75, thin]
      (-1.5,1.1) -- (5.3,1.1) -- (1.5,-1.1) -- (-5.3,-1.1) -- cycle;

    %, the line of solutions: the common intersection of the two planes ---
    \draw[ThemeAccent, very thick] (-4.3,0) -- (4.3,0);
    \node[ThemeAccent, anchor=west, inner sep=2pt] at (4.35,0) {$L$};

    %, the point singled out once a third equation is imposed ---
    \fill[ThemePurple] (1.6,0) circle (2.2pt);
    \node[ThemePurple, anchor=south, inner sep=4pt] at (1.6,0) {$c = 1$};

    \node[ThemeSecondary!75!black, anchor=south east] at (-4.6,1.5) {$R_1$};
    \node[ThemeGreen!62!black, anchor=south west] at (4.9,1.1) {$R_2$};
  \end{tikzpicture}
  \caption{Schematic picture of the motivating example. Each equation cuts out a
  plane in $\mathbb{R}^3$, and the solution set $L$ of the system is the line in
  which the two planes meet, one free parameter $c$, exactly as
  \eqref{eq:motivating_solution_set} records. Imposing a third independent
  equation adds a third plane, which meets $L$ in a single point.}
  \label{fig:two_planes}
\end{figure}

\begin{example}[One more equation removes the freedom]
Suppose we append a third equation $R_3: 3x + 2y + z = 0$ to the system above. Substituting the parametrised solution \eqref{eq:motivating_solution_set} into it gives
\[
3(2 - c) + 2(-1 - c) + c = 4 - 4c = 0,
\]
so that $c = 1$ is forced. The freedom collapses and the system now has the unique solution
\[
(x, y, z) = (1, -2, 1),
\]
which is exactly the member of the family \eqref{eq:motivating_solution_set} with $c = 1$. A system may therefore have no solution, exactly one solution, or infinitely many; deciding which of these occurs is one of the goals of this chapter.
\end{example}

\subsection{General Systems of Linear Equations}

\begin{definition}[Linear Equation]
\label{def:linear_equation}
A \newterm{linear equation} in $n$ variables $x_1, \dots, x_n$ over a field $K$ is an equation of the form
\[
a_1 x_1 + \dots + a_n x_n = b,
\]
where $a_1, \dots, a_n \in K$ are the \newterm{coefficients} and $b \in K$ is the \newterm{constant term}.
\begin{itemize}
    \item If $b = 0$, the equation is called \newterm{homogeneous}.
    \item If $b \ne 0$, the equation is called \newterm{inhomogeneous}.
\end{itemize}
\end{definition}

\begin{definition}[System of Linear Equations]
\label{def:system_of_linear_equations}
A general system of $m$ linear equations in $n$ variables over a field $K$ is written as
\begin{equation} \label{eq:general_system}
\begin{cases}
a_{1,1}x_1 + a_{1,2}x_2 + \dots + a_{1,n}x_n = b_1 \\
a_{2,1}x_1 + a_{2,2}x_2 + \dots + a_{2,n}x_n = b_2 \\
\vdots \\
a_{m,1}x_1 + a_{m,2}x_2 + \dots + a_{m,n}x_n = b_m
\end{cases}
\end{equation}
with $a_{i,j} \in K$ for $1 \le i \le m$, $1 \le j \le n$, and $b_i \in K$ for $1 \le i \le m$. Here $i$ is the row index and $j$ is the column index.
\end{definition}

\begin{definition}[Solution]
\label{def:solution}
A \newterm{solution} of the system is a tuple of scalars $s_1, \dots, s_n \in K$ that satisfies all $m$ equations simultaneously when we substitute $x_j = s_j$.
\end{definition}

\begin{definition}[Equivalence of Systems]
\label{def:equivalent_systems}
Two systems of linear equations in the same variables are called \newterm{equivalent} if they have exactly the same set of solutions.
\end{definition}

\subsection{Matrices}

To study systems of linear equations systematically, we strip away the variables and isolate the coefficients into a rectangular array.

\begin{definition}[Matrix]
\label{def:matrix}
An $m \times n$ \newterm{matrix} over $K$ is a rectangular array consisting of $m \cdot n$ elements of $K$, arranged in $m$ rows and $n$ columns.
\end{definition}

\begin{notation}
We denote a matrix $A$ by its entries,
\[
A = (a_{i,j})_{1 \le i \le m, \, 1 \le j \le n} \quad \text{or simply} \quad (a_{ij}),
\]
where $i$ denotes the row index and $j$ the column index.
\end{notation}

We define the set of all such matrices as
\begin{equation*}
    \M_{m \times n}(K) := \text{the set of all } m \times n \text{ matrices over } K.
\end{equation*}

\begin{definition}[Row and Column Vectors]
\label{def:row_and_column_vectors}
Special cases of matrices occur when one of the dimensions is $1$:
\begin{enumerate}[label=\textbf{(\alph*)}]
    \item An $n \times 1$ matrix is called a \newterm{column vector}. The set of all column vectors of length $n$ is denoted by $K^n := \M_{n \times 1}(K)$:
    \[
    \begin{pmatrix} x_1 \\ x_2 \\ \vdots \\ x_n \end{pmatrix} \in K^n.
    \]
    \item A $1 \times n$ matrix is called a \newterm{row vector}. The set of all row vectors of length $n$ is denoted by $K_{\text{row}}^n := \M_{1 \times n}(K)$:
    \[
    (x_1, x_2, \dots, x_n) \in K_{\text{row}}^n.
    \]
\end{enumerate}
\end{definition}

\subsection{Matrix-Vector Multiplication}

\begin{definition}[Matrix-Vector Multiplication]
\label{def:matrix_vector_multiplication}
Let $A \in \M_{m \times n}(K)$ and $x \in K^n$. We define the product $A \cdot x$ to be the column vector in $K^m$ obtained by pairing each row of $A$ with the vector $x$:
\[
\begin{pmatrix}
a_{1,1} & a_{1,2} & \dots & a_{1,n} \\
a_{2,1} & a_{2,2} & \dots & a_{2,n} \\
\vdots & \vdots & \ddots & \vdots \\
a_{m,1} & a_{m,2} & \dots & a_{m,n}
\end{pmatrix}
\begin{pmatrix} x_1 \\ x_2 \\ \vdots \\ x_n \end{pmatrix}
=
\begin{pmatrix}
a_{1,1}x_1 + \dots + a_{1,n}x_n \\
a_{2,1}x_1 + \dots + a_{2,n}x_n \\
\vdots \\
a_{m,1}x_1 + \dots + a_{m,n}x_n
\end{pmatrix} \in K^m.
\]
\end{definition}

With this notation, the whole system \eqref{eq:general_system} collapses into a single matrix equation. We write
\begin{equation} \label{eq:matrix_equation}
(S): \quad A x = b,
\end{equation}
where $A \in \M_{m \times n}(K)$ is the \newterm{coefficient matrix}, $x \in K^n$ is the vector of variables, and $b \in K^m$ is the \newterm{constant vector}.

\begin{notation}
\label{not:solution_set}
For a system $(S)$ as in \eqref{eq:matrix_equation} we write
\[
L(S) := \{(x_1, \dots, x_n) \mid x_i \in K, \; A x = b\} \subseteq K^n
\]
for its \newterm{set of solutions}. Two systems $(S)$ and $(S')$ in the same variables are equivalent (\cref{def:equivalent_systems}) precisely when $L(S) = L(S')$.
\end{notation}

\begin{example}[The motivating example in matrix form]
The system of the previous section can be rewritten first with indexed variables and then as a matrix equation:
\[
\begin{cases}
x + 3y + 4z = -1 \\
2x - y + z = 5
\end{cases}
\iff
\begin{cases}
x_1 + 3x_2 + 4x_3 = -1 \\
2x_1 - x_2 + x_3 = 5
\end{cases}
\iff
\underbrace{\begin{pmatrix} 1 & 3 & 4 \\ 2 & -1 & 1 \end{pmatrix}}_{2 \times 3}
\underbrace{\begin{pmatrix} x_1 \\ x_2 \\ x_3 \end{pmatrix}}_{3 \times 1}
=
\underbrace{\begin{pmatrix} -1 \\ 5 \end{pmatrix}}_{2 \times 1}.
\]
\end{example}

\begin{goals}
\label{goals:solution_set}
Our aim for the rest of this chapter is to study the set of solutions $L(S)$ of a system $A x = b$:
\begin{enumerate}[label=\textbf{(\alph*)}]
    \item Does the system have a solution at all, and if so, how many?
    \item How do we find \emph{all} of the solutions, systematically?
    \item Does $L(S)$ carry a structure of its own?
\end{enumerate}
\end{goals}

\subsection{Elementary Row Operations}

\begin{notation}
\label{not:augmented_matrix}
We do not need to write down the unknowns $x_1, \dots, x_n$ at every step: the coefficient matrix $A$ and the constant vector $b$ already carry all the information. We therefore record the system as its \newterm{augmented matrix} (or \newterm{extended matrix})
\[
(A \mid b) \in \M_{m \times (n+1)}(K),
\]
obtained by appending $b$ to $A$ as an extra column.
\end{notation}

\begin{example}[An augmented matrix]
The system
\[
\begin{cases}
2x_1 - \phantom{1}x_2 + 3x_3 + 2x_4 = 0 \\
\phantom{2}x_1 + 4x_2 \phantom{{}+ 3x_3} - \phantom{2}x_4 = 3 \\
2x_1 + 6x_2 - \phantom{3}x_3 + 5x_4 = 9
\end{cases}
\qquad \text{has the augmented matrix} \qquad
\begin{pmatrix}
2 & -1 & 3 & 2 & \big| & 0 \\
1 & 4 & 0 & -1 & \big| & 3 \\
2 & 6 & -1 & 5 & \big| & 9
\end{pmatrix}.
\]
\end{example}

The three operations below act on the equations of a system, and equally on the rows of its augmented matrix.

\begin{definition}[Elementary Row Operations]
\label{def:elementary_row_operations}
Let $R_i$ denote the $i$-th equation of a system, equivalently the $i$-th row of its augmented matrix. The three \newterm{elementary row operations} (EROs) are:
\begin{enumerate}[label=\textbf{(\alph*)}]
    \item \textbf{Type 1 (scaling).} Choose $\lambda \in K \setminus \{0\}$ and multiply row $i$ by $\lambda$. \\
    Notation: $\lambda \cdot R_i \to R_i$.
    \item \textbf{Type 2 (addition).} Choose $\lambda \in K$ and $1 \le i, j \le m$ with $i \ne j$, and replace row $j$ by row $j$ \emph{plus} $\lambda$ times row $i$. \\
    Notation: $R_j + \lambda \cdot R_i \to R_j$.
    \item \textbf{Type 3 (interchange).} Swap row $i$ and row $j$. \\
    Notation: $R_i \leftrightarrow R_j$.
\end{enumerate}
\end{definition}

\begin{definition}[Row Equivalence]
\label{def:row_equivalence}
Let $C$ and $C'$ be $r \times s$ matrices with entries in $K$. We say that $C'$ is \newterm{row equivalent} to $C$ if $C'$ can be obtained from $C$ by a finite sequence of elementary row operations,
\[
C = C_0 \longrightarrow C_1 \longrightarrow \dots \longrightarrow C_k = C'.
\]
\end{definition}

\begin{theorem}[Row Operations Preserve the Solution Set]
\label{thm:ero_preserve_solutions}
Let $(S): A x = b$ and $(S'): A' x = b'$ be two systems of $m$ linear equations in $n$ unknowns. Suppose the augmented matrix $(A' \mid b')$ is row equivalent to $(A \mid b)$. Then
\[
L(S') = L(S).
\]
\end{theorem}

\begin{proof}
We begin with three claims.

\begin{claim*}[Claim 1, One Operation Gives One Inclusion]
\label{claim:ero_one_inclusion}
If $(S')$ is obtained from $(S)$ by one elementary row operation, then $L(S) \subseteq L(S')$.
\end{claim*}

\emph{Proof of Claim 1.} Let $(x_1, \dots, x_n) \in L(S)$. We must show that $(x_1, \dots, x_n) \in L(S')$. We treat the three types separately.

\textbf{(a)} Operation $\lambda \cdot R_i \to R_i$ with $\lambda \ne 0$. All equations of $(S')$ coincide with those of $(S)$ except for equation $i$, which changes as
\[
a_{i,1}x_1 + \dots + a_{i,n}x_n = b_i
\qquad \longrightarrow \qquad
\lambda a_{i,1}x_1 + \dots + \lambda a_{i,n}x_n = \lambda b_i.
\]
If $(x_1, \dots, x_n)$ satisfies the left equation, then multiplying both sides by $\lambda$ shows that it satisfies the right one as well.

\textbf{(b)} Operation $R_j + \lambda \cdot R_i \to R_j$. Again all equations of $(S)$ and $(S')$ coincide except for equation $j$:
\begin{align*}
\text{row } j \text{ in } (S): \quad & a_{j,1}x_1 + \dots + a_{j,n}x_n = b_j \\
\longrightarrow \qquad \text{row } j \text{ in } (S'): \quad & \bigl(a_{j,1} + \lambda a_{i,1}\bigr)x_1 + \dots + \bigl(a_{j,n} + \lambda a_{i,n}\bigr)x_n = b_j + \lambda b_i
\end{align*}
If $(x_1, \dots, x_n)$ satisfies rows $i$ and $j$ of $(S)$, then adding $\lambda$ times the $i$-th equation to the $j$-th shows that it satisfies row $j$ of $(S')$ too.

\textbf{(c)} Operation $R_i \leftrightarrow R_j$. Here we clearly have
\[
(x_1, \dots, x_n) \in L(S) \implies (x_1, \dots, x_n) \in L(S'),
\]
because the order in which the equations of a system are listed has no effect whatsoever on its set of solutions.

This concludes the proof of Claim 1. \hfill $\diamond$

\begin{claim*}[Claim 2, Every Operation is Reversible]
\label{claim:ero_reversible}
If $(S')$ is obtained from $(S)$ by one elementary row operation, then $(S)$ can be obtained from $(S')$ by one elementary row operation.
\end{claim*}

\emph{Proof of Claim 2.} There are three cases, and in each of them we exhibit the inverse operation explicitly:
\begin{enumerate}[label=\textbf{(\alph*)}]
    \item If $(S) \xrightarrow{\;\lambda \cdot R_i \to R_i\;} (S')$ with $\lambda \ne 0$, then $(S') \xrightarrow{\;\frac{1}{\lambda} \cdot R_i \to R_i\;} (S)$.
    \item If $(S) \xrightarrow{\;R_j + \lambda \cdot R_i \to R_j\;} (S')$ with $i \ne j$, then $(S') \xrightarrow{\;R_j - \lambda \cdot R_i \to R_j\;} (S)$.
    \item If $(S) \xrightarrow{\;R_i \leftrightarrow R_j\;} (S')$, then $(S') \xrightarrow{\;R_j \leftrightarrow R_i\;} (S)$.
\end{enumerate}
Note that case \textbf{(a)} is where the restriction $\lambda \ne 0$ in \cref{def:elementary_row_operations} earns its keep: without it, $\frac{1}{\lambda}$ would not exist and the operation would not be reversible. This concludes the proof of Claim 2. \hfill $\diamond$

\begin{claim*}[Claim 3, One Operation Preserves the Solution Set]
\label{claim:ero_equality}
If $(S')$ is obtained from $(S)$ by one elementary row operation, then $L(S) = L(S')$.
\end{claim*}

\emph{Proof of Claim 3.} By Claim 1, we have $L(S) \subseteq L(S')$. By Claim 2, $(S)$ is obtained from $(S')$ by one elementary row operation, and hence Claim 1 applies again with the roles of $(S)$ and $(S')$ reversed, giving $L(S') \subseteq L(S)$. The two inclusions together give equality. This concludes the proof of Claim 3. \hfill $\diamond$

We are now in a position to prove the theorem. By assumption, there is a finite sequence of elementary row operations
\[
(S) = (S_0) \longrightarrow (S_1) \longrightarrow \dots \longrightarrow (S_k) = (S').
\]
Applying Claim 3 to each single step gives
\[
L(S) = L(S_0) = L(S_1) = \dots = L(S_k) = L(S'). \qedhere
\]
\end{proof}

\begin{exercise}[Converse of Solution Set Invariance under EROs]
\label{exc:converse_ero_solution_set}
We have seen in \cref{thm:ero_preserve_solutions} that if the augmented matrices $(A \mid b)$ and $(A' \mid b')$ of two linear systems of equations are row equivalent, then the two systems have the same solution set. Is the converse true: if two systems over a field $K$ have the same set of solutions, are their augmented matrices necessarily row equivalent? If yes, prove it; if not, provide a counterexample.
\exinfo{This exercise is Problem 4 of Problem Sheet 4, Linear Algebra I, Autumn Semester 2025.}
\exsol{sol:converse_ero_solution_set}
\end{exercise}

\subsection{Row-Reduced Matrices}

\Cref{thm:ero_preserve_solutions} tells us that we may row-reduce as much as we like without changing the answer. It remains to decide \emph{how far} the reduction can be pushed. The next two definitions describe the two successive normal forms we aim for.

\begin{definition}[Row-Reduced Matrix]
\label{def:row_reduced}
An $m \times n$ matrix $A$ is called \newterm{row-reduced} if the following two conditions hold:
\begin{enumerate}[label=\textbf{(\alph*)}]
    \item The first non-zero entry in each non-zero row of $A$ is $1$. This entry is called the \newterm{pivot} (or \newterm{leading term}) of that row.
    \item Each column of $A$ which contains the leading non-zero entry of some row has all its other entries equal to $0$.
\end{enumerate}
\end{definition}

\begin{example}[Row-reduced and not row-reduced]
\label{ex:row_reduced_examples}
Consider the following four matrices over $\mathbb{Q}$:
\[
A = \begin{pmatrix}
0 & 0 & 0 & 1 & 2 \\
0 & 1 & -3 & 0 & \frac{1}{2} \\
0 & 0 & 0 & 0 & 0
\end{pmatrix}
\qquad\qquad
B = \begin{pmatrix}
1 & 0 & 0 & 0 \\
0 & 1 & -1 & 0 \\
0 & 0 & 1 & 0
\end{pmatrix}
\]
\[
C = \begin{pmatrix}
0 & 2 & 1 \\
1 & 0 & -3 \\
0 & 0 & 0
\end{pmatrix}
\qquad\qquad
D = \begin{pmatrix}
0 & 0 & 1 & 1 \\
1 & 0 & 0 & 2 \\
0 & 1 & 0 & 5
\end{pmatrix}
\]
\begin{itemize}
    \item $A$ \emph{is} row-reduced. The pivots sit in columns $4$ and $2$, and each of those columns is zero apart from its pivot.
    \item $B$ is \emph{not} row-reduced. Its pivots sit in columns $1$, $2$ and $3$; but column $3$ contains the pivot of row $3$ and also the entry $-1$ in row $2$, so condition \textbf{(b)} fails.
    \item $C$ is \emph{not} row-reduced. The first non-zero entry of row $1$ is $2$ and not $1$, so condition \textbf{(a)} already fails.
    \item $D$ \emph{is} row-reduced. Its pivots sit in columns $3$, $1$ and $2$, and each of those columns is zero apart from its pivot. Note that the pivots do \emph{not} march to the right as we go down the rows, being row-reduced does not require this.
\end{itemize}
\end{example}

\begin{theorem}[Reduction to a Row-Reduced Matrix]
\label{thm:row_reduction_exists}
Every $m \times n$ matrix $A$ is row equivalent to a row-reduced matrix.
\end{theorem}

\begin{proof}
Let $A = (a_{ij})$ with $1 \le i \le m$, $1 \le j \le n$ and $a_{ij} \in K$. We proceed row by row.

\textbf{Row $\#1$.} If all entries in the first row of $A$ are $0$, then condition \textbf{(a)} of \cref{def:row_reduced} is satisfied for row $\#1$ and there is nothing to do. Otherwise, let $k$ be the smallest index $j$ for which $a_{1j} \ne 0$, that is
\[
a_{1j} = 0 \quad \forall \, 1 \le j < k, \qquad a_{1k} \ne 0,
\]
so that the first row has the shape $(0, \dots, 0, a_{1k}, \ast, \dots, \ast)$. Multiply row $\#1$ by $\frac{1}{a_{1k}}$; its leading entry is now $1$, and condition \textbf{(a)} holds for row $\#1$. Next, for each $2 \le i \le m$, add $(-a_{ik})$ times row $\#1$ to row $i$, that is, perform
\[
-a_{ik} \cdot R_1 + R_i \to R_i.
\]
Schematically:
\[
\begin{pmatrix}
0 & \dots & 0 & 1 & \dots \\
& \vdots & & \vdots & \\
a_{i1} & \dots & \dots & a_{ik} & \dots \\
& \vdots & & \vdots &
\end{pmatrix}
\xrightarrow{\; -a_{ik} \cdot R_1 + R_i \to R_i \;}
\begin{pmatrix}
0 & \dots & 0 & 1 & \dots \\
& \vdots & & \vdots & \\
a_{i1} & \dots & a_{i,k-1} & 0 & \dots \\
& \vdots & & \vdots &
\end{pmatrix}
\]
After these operations, column $k$ is zero apart from the $1$ in row $\#1$.

\textbf{Row $\#2$.} If all entries in row $\#2$ are $0$, we do nothing to this row. If some entry in row $\#2$ is non-zero, we multiply row $\#2$ by a scalar so that its leading term becomes $1$. Recall that row $\#1$ has its leading non-zero entry in column $k$. Since we have just made all other entries of column $k$ vanish, the leading non-zero term of row $\#2$ can \emph{not} lie in column $k$: it must appear in some other column, say column $k'$ with $k' \ne k$. We now add suitable multiples of row $\#2$ to all the other rows until we obtain a matrix in which all entries of column $k'$ are $0$, except the one in row $\#2$. The two possible configurations are
\[
\begin{pmatrix}
0 & \dots & 0 & \dots & 0 & 1 & \ast & \dots & \ast \\
0 & \dots & 0 & 1 & \ast \dots \ast & 0 & \ast & \dots & \ast \\
& & 0 & & & 0 & & & \\
& & \vdots & & & \vdots & & &
\end{pmatrix}
\quad \text{or} \quad
\begin{pmatrix}
0 & \dots & 0 & 1 & \ast \dots \ast & 0 & \ast & \dots & \ast \\
0 & \dots & 0 & \dots & 0 & 1 & \ast & \dots & \ast \\
& & 0 & & & 0 & & & \\
& & \vdots & & & \vdots & & &
\end{pmatrix}
\]
according to whether $k' < k$ or $k < k'$.

\begin{importantremark}
\label{rem:row_reduction_does_not_undo}
In the operations just performed using row $\#2$ we did \emph{not} change any of the entries of row $\#1$ lying in columns $1, \dots, k$; in particular, nothing was changed in column $k$. Also, in case row $\#1$ was zero, the operations using row $\#2$ do not affect row $\#1$ at all. This is what guarantees that the work already done is never undone by the later steps.
\end{importantremark}

We continue now to row $\#3$ and repeat the same procedure, and so on. After a finite number of elementary row operations we arrive at a row-reduced matrix.
\end{proof}

\subsection{Row-Reduced Echelon Matrices}

The matrix $D$ of the example above is row-reduced, and yet its pivots appear in the disorderly column order $3, 1, 2$. Sorting the rows costs nothing and yields a genuine normal form.

\begin{definition}[Row-Reduced Echelon Matrix]
\label{def:row_reduced_echelon}
An $m \times n$ matrix $A$ is called \newterm{row-reduced echelon} if the following hold:
\begin{enumerate}[label=\textbf{(\alph*)}]
    \item $A$ is row-reduced.
    \item Every row of $A$ which has all its entries $0$ appears below every row which has a non-zero entry.
    \item If rows $1, \dots, r$ are the non-zero rows of $A$, and if the leading non-zero entry of row $i$ occurs in column $k_i$ for $i = 1, \dots, r$, then
    \[
    k_1 < k_2 < \dots < k_r.
    \]
\end{enumerate}
\end{definition}

Condition \textbf{(c)} is what produces the characteristic staircase shape shown in \cref{fig:echelon_staircase}.

\begin{figure}[ht!]
  \centering
  \begin{tikzpicture}[font=\small, x=1cm, y=0.66cm]
    % Cell (column c, row r) is centred at (c-0.5, -r+0.5).
    % Pivots sit in columns 2, 4, 6 of rows 1, 2, 3; rows 4 and 5 are zero.
    %, the region that the echelon conditions force to be zero ---
    \fill[gray!14]
      (0,0) -- (1,0) -- (1,-1) -- (3,-1) -- (3,-2) -- (5,-2) -- (5,-3)
      -- (7,-3) -- (7,-5) -- (0,-5) -- cycle;
    %, boundary between the non-zero rows and the zero rows ---
    \draw[ThemeGreen!55, dashed] (0,-3) -- (5,-3);
    %, zero entries ---
    \foreach \c/\r in {1/1, 1/2, 1/3, 1/4, 1/5,
                       2/2, 2/3, 3/2, 3/3, 4/1, 4/3, 5/3, 6/1, 6/2,
                       2/4, 3/4, 4/4, 5/4, 6/4, 7/4,
                       2/5, 3/5, 4/5, 5/5, 6/5, 7/5}
      \node[gray] at (\c-0.5,-\r+0.5) {$0$};
    %, entries in the pivot-free columns ---
    \foreach \c/\r in {3/1, 5/1, 7/1, 5/2, 7/2, 7/3}
      \node[ThemeSecondary] at (\c-0.5,-\r+0.5) {$\ast$};
    %, the pivots ---
    \foreach \c/\r in {2/1, 4/2, 6/3} {
      \node[ThemeAccent] at (\c-0.5,-\r+0.5) {$1$};
      \draw[ThemeAccent, thick] (\c-0.5,-\r+0.5) circle (0.3);
    }
    %, the staircase itself, drawn last so that it stays visible ---
    \draw[ThemeGreen, very thick]
      (1,0) -- (1,-1) -- (3,-1) -- (3,-2) -- (5,-2) -- (5,-3) -- (7,-3);
    %, delimiters ---
    \draw[thick] (-0.15,0.15) -- (-0.4,0.15) -- (-0.4,-5.15) -- (-0.15,-5.15);
    \draw[thick] (7.15,0.15) -- (7.4,0.15) -- (7.4,-5.15) -- (7.15,-5.15);
    %, row labels ---
    \node[anchor=east, gray] at (-0.7,-0.5) {$1$};
    \node[anchor=east, gray] at (-0.7,-2.5) {$r$};
    \node[anchor=east, gray] at (-0.7,-3.5) {$r+1$};
    \node[anchor=east, gray] at (-0.7,-4.5) {$m$};
    %, pivot column labels ---
    \foreach \c/\lab in {2/{k_1}, 4/{k_2}, 6/{k_3}}
      \node[ThemeAccent, anchor=south] at (\c-0.5,0.25) {$\lab$};
  \end{tikzpicture}
  \caption{A row-reduced echelon matrix. The circled pivots are $1$; everything
  below the green staircase is $0$, and each pivot column is $0$ apart from its
  pivot. Because $k_1 < k_2 < k_3$, the pivots move strictly to the right as we
  descend, and the $m - r$ zero rows are pushed to the bottom. The columns
  carrying a $\ast$ are the ones without a pivot; they will turn into the free
  variables when the matrix is read back as a system.}
  \label{fig:echelon_staircase}
\end{figure}

\begin{theorem}[Reduction to Row-Reduced Echelon Form]
\label{thm:echelon_form_exists}
Every $m \times n$ matrix $A$ is row equivalent to a row-reduced echelon matrix.
\end{theorem}

\begin{proof}
By \cref{thm:row_reduction_exists} we already know that $A$ is row equivalent to a row-reduced matrix $A'$. Performing a finite number of row interchanges on $A'$ (an operation of Type 3, which preserves row equivalence), we may list the non-zero rows in increasing order of their pivot columns and move all zero rows to the bottom. The result satisfies conditions \textbf{(b)} and \textbf{(c)} of \cref{def:row_reduced_echelon}, and interchanging rows does not disturb condition \textbf{(a)}. Hence, we arrive at a row-reduced echelon matrix.
\end{proof}

\begin{remark}
\label{rem:why_echelon_helps}
Why is this useful for solving systems of equations? Let $A x = b$ be a system of linear equations and consider its augmented matrix $(A \mid b)$. We may perform on $(A \mid b)$ the same elementary operations that we perform on $A$, and so obtain a row-equivalent matrix $(A' \mid b')$ in which $A'$ is a row-reduced echelon matrix. By \cref{thm:ero_preserve_solutions}, the systems $A x = b$ and $A' x = b'$ have the same solutions, but $A' x = b'$ is far easier to solve, since each pivot expresses one variable directly in terms of the pivot-free ones.
\end{remark}

\subsection{Worked Examples}

\begin{example}[Solving a system by full reduction]
\label{ex:full_reduction}
Consider the system over $\mathbb{R}$ whose augmented matrix is
\[
\begin{pmatrix}
0 & -9 & 3 & 4 & \big| & 9 \\
1 & 4 & 0 & -1 & \big| & 5 \\
2 & 6 & -1 & 5 & \big| & -5
\end{pmatrix}.
\]
We first normalise the leading entry of the first row, and then clear its column:
\[
\xrightarrow{-\frac{1}{9} \cdot R_1 \to R_1}
\begin{pmatrix}
0 & 1 & -\frac{1}{3} & -\frac{4}{9} & \big| & -1 \\
1 & 4 & 0 & -1 & \big| & 5 \\
2 & 6 & -1 & 5 & \big| & -5
\end{pmatrix}
\xrightarrow{\begin{subarray}{l} R_2 - 4 \cdot R_1 \to R_2 \\ R_3 - 6 \cdot R_1 \to R_3 \end{subarray}}
\begin{pmatrix}
0 & 1 & -\frac{1}{3} & -\frac{4}{9} & \big| & -1 \\
1 & 0 & \frac{4}{3} & \frac{7}{9} & \big| & 9 \\
2 & 0 & 1 & \frac{23}{3} & \big| & 1
\end{pmatrix}
\]
Next, we clear the first column below row $2$ and normalise the third row:
\[
\xrightarrow{R_3 - 2 \cdot R_2 \to R_3}
\begin{pmatrix}
0 & 1 & -\frac{1}{3} & -\frac{4}{9} & \big| & -1 \\
1 & 0 & \frac{4}{3} & \frac{7}{9} & \big| & 9 \\
0 & 0 & -\frac{5}{3} & \frac{55}{9} & \big| & -17
\end{pmatrix}
\xrightarrow{-\frac{3}{5} \cdot R_3 \to R_3}
\begin{pmatrix}
0 & 1 & -\frac{1}{3} & -\frac{4}{9} & \big| & -1 \\
1 & 0 & \frac{4}{3} & \frac{7}{9} & \big| & 9 \\
0 & 0 & 1 & -\frac{11}{3} & \big| & \frac{51}{5}
\end{pmatrix}
\]
Now we clear the third column upwards, and finally sort the rows:
\[
\xrightarrow{\begin{subarray}{l} R_2 - \frac{4}{3} \cdot R_3 \to R_2 \\R_1 + \frac{1}{3} \cdot R_3 \to R_1 \end{subarray}}
\begin{pmatrix}
0 & 1 & 0 & -\frac{5}{3} & \big| & \frac{12}{5} \\
1 & 0 & 0 & \frac{17}{3} & \big| & -\frac{23}{5} \\
0 & 0 & 1 & -\frac{11}{3} & \big| & \frac{51}{5}
\end{pmatrix}
\xrightarrow{R_1 \leftrightarrow R_2}
\begin{pmatrix}
1 & 0 & 0 & \frac{17}{3} & \big| & -\frac{23}{5} \\
0 & 1 & 0 & -\frac{5}{3} & \big| & \frac{12}{5} \\
0 & 0 & 1 & -\frac{11}{3} & \big| & \frac{51}{5}
\end{pmatrix}
\]
The last matrix is in row-reduced echelon form. Reading it back as a system gives
\[
\begin{cases}
x_1 \phantom{{}+ x_2 + x_3} + \frac{17}{3}x_4 = -\frac{23}{5} \\
\phantom{x_1 + {}} x_2 \phantom{{}+ x_3} - \frac{5}{3}x_4 = \phantom{-}\frac{12}{5} \\
\phantom{x_1 + x_2 + {}} x_3 - \frac{11}{3}x_4 = \phantom{-}\frac{51}{5}
\end{cases}
\]
Column $4$ carries no pivot, so $x_4$ is free. Take $x_4 := a \in \mathbb{R}$ arbitrary. Then
\[
x_3 = \frac{11}{3}a + \frac{51}{5}, \qquad
x_2 = \frac{5}{3}a + \frac{12}{5}, \qquad
x_1 = -\frac{17}{3}a - \frac{23}{5},
\]
and these are \emph{all} the solutions of the original system, by \cref{thm:ero_preserve_solutions}.
\end{example}

\begin{exercise}[Gaussian Elimination on an Underdetermined System]
\label{exc:gauss_elimination_4vars}
Using Gaussian elimination, find the set of all solutions in $\mathbb{R}^4$ to the following system of linear equations over $\mathbb{R}$:
\[
\begin{cases}
\phantom{3}x + 2y + 2z + \phantom{5}w = -1 \\
3x + 6y + 2z + 5w = \phantom{-}1
\end{cases}
\]
\exinfo{This exercise is Problem 1 of Problem Sheet 3, Linear Algebra I, Autumn Semester 2025.}
\exsol{sol:gauss_elimination_4vars}
\end{exercise}

\begin{exercise}[Simultaneous Row Reduction for Multiple Right-Hand Sides]
\label{exc:simultaneous_row_reduction}
Consider the following two systems of linear equations over $\mathbb{R}$:
\[
(S_1): \begin{cases}
\phantom{2}x + \phantom{2}y + \phantom{2}z = 6 \\
\phantom{2}x + 2y + 2z = 11 \\
2x + 3y - 4z = 3
\end{cases}
\qquad\text{and}\qquad
(S_2): \begin{cases}
\phantom{2}x + \phantom{2}y + \phantom{2}z = 7 \\
\phantom{2}x + 2y + 2z = 10 \\
2x + 3y - 4z = 3
\end{cases}
\]
Solve both systems simultaneously by row-reducing the single $3 \times 5$ augmented matrix $(A \mid b_1 \mid b_2)$.
\exinfo{This exercise is Problem 3 of Problem Sheet 4, Linear Algebra I, Autumn Semester 2025.}
\exsol{sol:simultaneous_row_reduction}
\end{exercise}

\begin{example}[Conditions for consistency]
\label{ex:consistency_conditions}
Let us now see how the solutions behave for a general choice of the constants $b_1, b_2, b_3 \in \mathbb{R}$ in the system
\[
\begin{pmatrix}
1 & -2 & 1 \\
2 & 1 & 1 \\
0 & 5 & -1
\end{pmatrix}
\begin{pmatrix} x_1 \\ x_2 \\ x_3 \end{pmatrix}
= \begin{pmatrix} b_1 \\ b_2 \\ b_3 \end{pmatrix}.
\]
We construct the augmented matrix and reduce it:
\[
\begin{pmatrix}
1 & -2 & 1 & \big| & b_1 \\
2 & 1 & 1 & \big| & b_2 \\
0 & 5 & -1 & \big| & b_3
\end{pmatrix}
\xrightarrow{R_2 - 2 \cdot R_1 \to R_2}
\begin{pmatrix}
1 & -2 & 1 & \big| & b_1 \\
0 & 5 & -1 & \big| & b_2 - 2b_1 \\
0 & 5 & -1 & \big| & b_3
\end{pmatrix}
\]
\[
\xrightarrow{R_3 - R_2 \to R_3}
\begin{pmatrix}
1 & -2 & 1 & \big| & b_1 \\
0 & 5 & -1 & \big| & b_2 - 2b_1 \\
0 & 0 & 0 & \big| & b_3 - b_2 + 2b_1
\end{pmatrix}
\xrightarrow{\frac{1}{5} \cdot R_2 \to R_2}
\begin{pmatrix}
1 & -2 & 1 & \big| & b_1 \\
0 & 1 & -\frac{1}{5} & \big| & \frac{1}{5}(b_2 - 2b_1) \\
0 & 0 & 0 & \big| & b_3 - b_2 + 2b_1
\end{pmatrix}
\]
\[
\xrightarrow{R_1 + 2 \cdot R_2 \to R_1}
\begin{pmatrix}
1 & 0 & \frac{3}{5} & \big| & \frac{1}{5}(b_1 + 2b_2) \\
0 & 1 & -\frac{1}{5} & \big| & \frac{1}{5}(b_2 - 2b_1) \\
0 & 0 & 0 & \big| & b_3 - b_2 + 2b_1
\end{pmatrix}
\]
The last matrix is in row-reduced echelon form, and the equivalent system reads
\[
\begin{cases}
x_1 \phantom{{}+ x_2} + \frac{3}{5}x_3 = \frac{1}{5}(b_1 + 2b_2) \\
\phantom{x_1 + {}} x_2 - \frac{1}{5}x_3 = \frac{1}{5}(b_2 - 2b_1) \\
\phantom{x_1 + x_2 + \frac{3}{5}x_3 \; } 0 = b_3 - b_2 + 2b_1
\end{cases}
\]
Two cases arise from the last row.
\begin{itemize}
    \item If $b_3 - b_2 + 2b_1 \ne 0$, the third equation reads $0 = \text{non-zero}$, and the system has \emph{no} solutions.
    \item If $b_3 - b_2 + 2b_1 = 0$, the third equation is vacuous and the system is equivalent to
    \[
    \begin{cases}
    x_1 + \frac{3}{5}x_3 = \frac{1}{5}(b_1 + 2b_2) \\
    x_2 - \frac{1}{5}x_3 = \frac{1}{5}(b_2 - 2b_1)
    \end{cases}
    \]
    Column $3$ carries no pivot, so $x_3$ is free. Take $x_3 := c \in \mathbb{R}$ arbitrary; then
    \[
    x_2 = \frac{1}{5}c + \frac{1}{5}(b_2 - 2b_1), \qquad
    x_1 = -\frac{3}{5}c + \frac{1}{5}(b_1 + 2b_2).
    \]
\end{itemize}
So, the constants $b_1, b_2, b_3$ cannot be chosen freely: the system is consistent if and only if $b_3 - b_2 + 2b_1 = 0$, and when it is, it has infinitely many solutions.
\end{example}

\begin{exercise}[Solvability and Solutions of a Two-Parameter System]
\label{exc:two_parameter_linear_system}
For which values of the parameters $a \in \mathbb{R}$ and $b \in \mathbb{R}$ does the following system of linear equations admit at least one solution $(x, y, z) \in \mathbb{R}^3$?
\[
\begin{cases}
\phantom{-}x + 2y - 3z = 10 \\
\phantom{-}3x + 7y - 9z = b \\
-3x - 5y + az = b \\
\phantom{-}2x + 4y - 3z = 5
\end{cases}
\]
For all such pairs $(a, b)$, determine the complete solution set.
\exinfo{This exercise is Problem 2 of Problem Sheet 4, Linear Algebra I, Autumn Semester 2025.}
\exsol{sol:two_parameter_linear_system}
\end{exercise}

\begin{exercise}[Parametric Solvability of a Linear System]
\label{exc:parametric_linear_system_solvability}
Fix constants $b_1, b_2, b_3 \in \mathbb{R}$. Consider the system of linear equations
\[
\begin{pmatrix}
0 & 7 & -2 \\
-1 & 2 & -\frac{1}{2} \\
4 & -1 & 0
\end{pmatrix}
\begin{pmatrix} x \\ y \\ z \end{pmatrix}
= \begin{pmatrix} b_1 \\ b_2 \\ b_3 \end{pmatrix}.
\]
Determine necessary and sufficient conditions on $b_1, b_2, b_3$ for this system to possess at least one solution, and describe the solution set $S \subseteq \mathbb{R}^3$ when solutions exist.
\exinfo{This exercise is Problem 4 of Problem Sheet 3, Linear Algebra I, Autumn Semester 2025.}
\exsol{sol:parametric_linear_system_solvability}
\end{exercise}

\begin{example}[A system with infinitely many solutions]
\label{ex:infinitely_many_solutions}
As a final illustration, consider
\begin{equation*}
\begin{cases}
x_1 - 2x_2 + \phantom{2}x_3 + \phantom{3}x_4 = 1 \\
2x_1 - 4x_2 + \phantom{2}x_3 + 3x_4 = 2 \\
x_1 - 2x_2 + 2x_3 - \phantom{3}x_4 = 1
\end{cases}
\end{equation*}
We write the augmented matrix and apply EROs:
\[
\begin{pmatrix}
1 & -2 & 1 & 1 & \big| & 1 \\
2 & -4 & 1 & 3 & \big| & 2 \\
1 & -2 & 2 & -1 & \big| & 1
\end{pmatrix}
\xrightarrow{\begin{subarray}{l} R_2 - 2 \cdot R_1 \to R_2 \\ R_3 - R_1 \to R_3 \end{subarray}}
\begin{pmatrix}
1 & -2 & 1 & 1 & \big| & 1 \\
0 & 0 & -1 & 1 & \big| & 0 \\
0 & 0 & 1 & -2 & \big| & 0
\end{pmatrix}
\]
\[
\xrightarrow{R_3 + R_2 \to R_3}
\begin{pmatrix}
1 & -2 & 1 & 1 & \big| & 1 \\
0 & 0 & -1 & 1 & \big| & 0 \\
0 & 0 & 0 & -1 & \big| & 0
\end{pmatrix}
\]
From the third row, $-x_4 = 0$, and therefore, $x_4 = 0$. From the second row, $-x_3 + x_4 = 0$, and therefore, $x_3 = 0$. Substituting into the first row gives $x_1 - 2x_2 = 1$, that is, $x_1 = 1 + 2x_2$.

Here $x_2$ is the free variable. Let $x_2 := \alpha \in \mathbb{R}$. The general solution is
\[
x = \begin{pmatrix} x_1 \\ x_2 \\ x_3 \\ x_4 \end{pmatrix} = \begin{pmatrix} 1 + 2\alpha \\ \alpha \\ 0 \\ 0 \end{pmatrix} = \begin{pmatrix} 1 \\ 0 \\ 0 \\ 0 \end{pmatrix} + \alpha \begin{pmatrix} 2 \\ 1 \\ 0 \\ 0 \end{pmatrix}.
\]
Observe the shape of the answer: one particular solution plus an arbitrary multiple of a fixed vector. That the solution set of a system should always look like this is precisely the kind of structure our \textbf{(c)} question asked about, in the list of goals on \cpageref{goals:solution_set}, and we will return to it once we have the language of vector spaces. The promise is kept in ch. 17d, where \cref{cor:non_homogeneous_solutions} says exactly this for an arbitrary consistent system.
\end{example}

\begin{exercise}[Inconsistency in \texorpdfstring{$2 \times 2$}{2x2} Systems and Linear Dependence]
\label{exc:inconsistent_2x2_system}
Let $A \in M_{2 \times 2}(\mathbb{R})$ be a $2 \times 2$ real matrix and let $b \in \mathbb{R}^2$ be a vector such that the system of equations $Ax = b$ has no solutions. Show that one row of the matrix $A$ must be a scalar multiple of the other.
\exinfo{This exercise is Problem 5 of Problem Sheet 4, Linear Algebra I, Autumn Semester 2025.}
\exsol{sol:inconsistent_2x2_system}
\end{exercise}

\begin{ainote}
The chapter establishes exactly one thing the method needs,
\cref{thm:ero_preserve_solutions}, that row operations change nothing about the
solution set. Everything else here is bookkeeping in its service, and its proof
through the three `claim*' blocks is the model for the reversibility arguments
used again in ch. 17e.

What it raises and cannot answer is worth noting, because the three worked
examples exhibit three different behaviours and none is explained. Why those are
the only shapes a solution set can take, why the number of free parameters is $n$
minus the number of pivots, and why that number does not depend on which sequence
of row operations was chosen, all wait for the rank of a matrix, in chs. 13 and
17d.

Two conventions are worth fixing in the mind before the computations begin.
Prof. Biran writes the base field $F$, where this book writes $K$ throughout, as
all the later chapters do. And in the three elementary row operations
$c \cdot R_i \to R_i$, $R_j + c \cdot R_i \to R_j$ and
$R_i \leftrightarrow R_j$, the row that changes always stands \emph{after} the
arrow: the second of them alters $R_j$ and leaves $R_i$ exactly as it was, which
is the reverse of what reading the expression left to right suggests.
\end{ainote}

\section{Solutions}\label{sec:solutions_linear_equations}

% Solution by Gemini 3.7 Flash (High)
\begin{exercisesolution}[Solution to \cref{exc:converse_ero_solution_set}]
\label{sol:converse_ero_solution_set}
\exinfo{Solution by Gemini 3.7 Flash (High). Examines row equivalence vs solution set equality for inconsistent and non-square systems.}
The converse is \textbf{false} in general. Having identical solution sets does not imply that the augmented matrices are row equivalent.

\textbf{Counterexample 1 (Inconsistent Systems):}
Consider the two linear systems over $\mathbb{R}$ in two variables $(x_1, x_2)$:
\[
(S_1): \begin{cases} x_1 + 0x_2 = 0 \\ 0x_1 + 0x_2 = 1 \end{cases} \qquad\text{and}\qquad (S_2): \begin{cases} 0x_1 + x_2 = 0 \\ 0x_1 + 0x_2 = 1 \end{cases}
\]
Both systems are inconsistent, so their solution sets are both empty: $L(S_1) = \emptyset = L(S_2)$.
Their augmented matrices in row-reduced echelon form are:
\[
(A_1 \mid b_1) = \begin{pmatrix} 1 & 0 & \big| & 0 \\ 0 & 0 & \big| & 1 \end{pmatrix} \qquad\text{and}\qquad (A_2 \mid b_2) = \begin{pmatrix} 0 & 1 & \big| & 0 \\ 0 & 0 & \big| & 1 \end{pmatrix}.
\]
These two matrices are not row equivalent because the only non-zero row in the coefficient part of $(A_1 \mid b_1)$ is $(1, 0)$, whereas in $(A_2 \mid b_2)$ it is $(0, 1)$. Since both matrices are in row-reduced echelon form (\cref{def:row_reduced_echelon}), they represent different row spaces and cannot be transformed into one another by elementary row operations.

\textbf{Counterexample 2 (Different Number of Equations):}
The system $(x_1 = 0)$ with augmented matrix $(1 \mid 0) \in \M_{1 \times 2}(\mathbb{R})$ and the system $\begin{cases} x_1 = 0 \\ x_1 = 0 \end{cases}$ with augmented matrix $\begin{pmatrix} 1 & 0 \\ 1 & 0 \end{pmatrix} \in \M_{2 \times 2}(\mathbb{R})$ both have the unique solution $x_1 = 0$. However, matrices of different dimensions cannot be row equivalent.

\begin{remark}
If one restricts attention to \emph{consistent} linear systems with the same number of variables and equations whose coefficient matrices have the same dimensions, then the converse does hold: any consistent system has a unique reduced row echelon form completely determined by the affine subspace of solutions.
\end{remark}
\end{exercisesolution}

% Solution by Gemini 3.7 Flash (High)
\begin{exercisesolution}[Solution to \cref{exc:gauss_elimination_4vars}]
\label{sol:gauss_elimination_4vars}
\exinfo{Solution by Gemini 3.7 Flash (High). Fully works out Gaussian elimination into reduced row echelon form with 2 free variables.}
We write the augmented matrix of the system and apply elementary row operations:
\[
\begin{pmatrix}
1 & 2 & 2 & 1 & \big| & -1 \\
3 & 6 & 2 & 5 & \big| & 1
\end{pmatrix}
\xrightarrow{-3 \cdot E_1 + E_2 \to E_2}
\begin{pmatrix}
1 & 2 & 2 & 1 & \big| & -1 \\
0 & 0 & -4 & 2 & \big| & 4
\end{pmatrix}
\]
Divide the second row by $-4$:
\[
\xrightarrow{-\frac{1}{4} \cdot E_2 \to E_2}
\begin{pmatrix}
1 & 2 & 2 & 1 & \big| & -1 \\
0 & 0 & 1 & -\frac{1}{2} & \big| & -1
\end{pmatrix}
\]
Now clear the entry above the second pivot by $-2 \cdot E_2 + E_1 \to E_1$:
\[
\xrightarrow{-2 \cdot E_2 + E_1 \to E_1}
\begin{pmatrix}
1 & 2 & 0 & 2 & \big| & 1 \\
0 & 0 & 1 & -\frac{1}{2} & \big| & -1
\end{pmatrix}
\]
The matrix is now in row-reduced echelon form. 
\begin{itemize}
    \item \textbf{Pivot variables:} $x$ (column 1) and $z$ (column 3).
    \item \textbf{Free variables:} $y$ (column 2) and $w$ (column 4).
\end{itemize}
Let $y = s \in \mathbb{R}$ and $w = t \in \mathbb{R}$ be arbitrary parameters. Expressing the pivot variables in terms of the free parameters gives:
\[
\begin{cases}
x = 1 - 2s - 2t, \\
z = -1 + \frac{1}{2}t.
\end{cases}
\]
Therefore, the complete solution set in $\mathbb{R}^4$ is:
\[
L = \left\{ \begin{pmatrix} x \\ y \\ z \\ w \end{pmatrix} = \begin{pmatrix} 1 - 2s - 2t \\ s \\ -1 + \frac{1}{2}t \\ t \end{pmatrix} = \begin{pmatrix} 1 \\ 0 \\ -1 \\ 0 \end{pmatrix} + s \begin{pmatrix} -2 \\ 1 \\ 0 \\ 0 \end{pmatrix} + t \begin{pmatrix} -2 \\ 0 \\ \frac{1}{2} \\ 1 \end{pmatrix} \;\middle|\; s, t \in \mathbb{R} \right\}.
\]
\end{exercisesolution}

% Solution by Gemini 3.7 Flash (High)
\begin{exercisesolution}[Solution to \cref{exc:simultaneous_row_reduction}]
\label{sol:simultaneous_row_reduction}
\exinfo{Solution by Gemini 3.7 Flash (High). Applies simultaneous Gaussian elimination on a dual augmented matrix.}
Since both systems $(S_1)$ and $(S_2)$ share the same coefficient matrix $A$, we form the $3 \times 5$ augmented matrix $(A \mid b_1 \mid b_2)$ and perform row reduction:
\[
\begin{pmatrix}
1 & 1 & 1 & \big| & 6 & 7 \\
1 & 2 & 2 & \big| & 11 & 10 \\
2 & 3 & -4 & \big| & 3 & 3
\end{pmatrix}
\xrightarrow{\begin{subarray}{l} -1 \cdot E_1 + E_2 \to E_2 \\ -2 \cdot E_1 + E_3 \to E_3 \end{subarray}}
\begin{pmatrix}
1 & 1 & 1 & \big| & 6 & 7 \\
0 & 1 & 1 & \big| & 5 & 3 \\
0 & 1 & -6 & \big| & -9 & -11
\end{pmatrix}
\]
Eliminate the second column in row 3:
\[
\xrightarrow{-1 \cdot E_2 + E_3 \to E_3}
\begin{pmatrix}
1 & 1 & 1 & \big| & 6 & 7 \\
0 & 1 & 1 & \big| & 5 & 3 \\
0 & 0 & -7 & \big| & -14 & -14
\end{pmatrix}
\xrightarrow{-\frac{1}{7} \cdot E_3 \to E_3}
\begin{pmatrix}
1 & 1 & 1 & \big| & 6 & 7 \\
0 & 1 & 1 & \big| & 5 & 3 \\
0 & 0 & 1 & \big| & 2 & 2
\end{pmatrix}
\]
Now clear the entries above the third pivot:
\[
\xrightarrow{\begin{subarray}{l} -1 \cdot E_3 + E_2 \to E_2 \\ -1 \cdot E_3 + E_1 \to E_1 \end{subarray}}
\begin{pmatrix}
1 & 1 & 0 & \big| & 4 & 5 \\
0 & 1 & 0 & \big| & 3 & 1 \\
0 & 0 & 1 & \big| & 2 & 2
\end{pmatrix}
\xrightarrow{-1 \cdot E_2 + E_1 \to E_1}
\begin{pmatrix}
1 & 0 & 0 & \big| & 1 & 4 \\
0 & 1 & 0 & \big| & 3 & 1 \\
0 & 0 & 1 & \big| & 2 & 2
\end{pmatrix}
\]
Reading off the solutions from the respective right-hand columns:
\begin{itemize}
    \item For $(S_1)$, the unique solution is $(x, y, z) = (1, 3, 2)$.
    \item For $(S_2)$, the unique solution is $(x, y, z) = (4, 1, 2)$.
\end{itemize}
\end{exercisesolution}

% Solution by Gemini 3.7 Flash (High)
\begin{exercisesolution}[Solution to \cref{exc:two_parameter_linear_system}]
\label{sol:two_parameter_linear_system}
\exinfo{Solution by Gemini 3.7 Flash (High). Analyzes consistency parameters $a, b$ and determines the unique solution.}
We write the augmented matrix of the system and apply Gaussian elimination:
\[
\begin{pmatrix}
1 & 2 & -3 & \big| & 10 \\
3 & 7 & -9 & \big| & b \\
-3 & -5 & a & \big| & b \\
2 & 4 & -3 & \big| & 5
\end{pmatrix}
\xrightarrow{\begin{subarray}{l} -3 \cdot E_1 + E_2 \to E_2 \\ 3 \cdot E_1 + E_3 \to E_3 \\ -2 \cdot E_1 + E_4 \to E_4 \end{subarray}}
\begin{pmatrix}
1 & 2 & -3 & \big| & 10 \\
0 & 1 & 0 & \big| & b - 30 \\
0 & 1 & a - 9 & \big| & b + 30 \\
0 & 0 & 3 & \big| & -15
\end{pmatrix}
\]
From row 4, we have $3z = -15 \implies z = -5$.
Next, perform $-1 \cdot E_2 + E_3 \to E_3$:
\[
\xrightarrow{-1 \cdot E_2 + E_3 \to E_3}
\begin{pmatrix}
1 & 2 & -3 & \big| & 10 \\
0 & 1 & 0 & \big| & b - 30 \\
0 & 0 & a - 9 & \big| & 60 \\
0 & 0 & 3 & \big| & -15
\end{pmatrix}
\]
The third row represents the equation $(a - 9)z = 60$. Substituting $z = -5$ gives:
\[
(a - 9)(-5) = 60 \implies a - 9 = -12 \implies a = -3.
\]
Therefore:
\begin{itemize}
    \item If $a \ne -3$, the third and fourth rows are contradictory, and the system has \textbf{no solutions}.
    \item If $a = -3$, the system is consistent for \textbf{every} $b \in \mathbb{R}$.
\end{itemize}
When $a = -3$, we determine the unique solution:
\begin{itemize}
    \item $z = -5$,
    \item $y = b - 30$,
    \item From row 1: $x + 2(b - 30) - 3(-5) = 10 \implies x + 2b - 60 + 15 = 10 \implies x = 55 - 2b$.
\end{itemize}
In conclusion, the system admits a solution if and only if $a = -3$ (with $b \in \mathbb{R}$ arbitrary), and the unique solution is $(x, y, z) = (55 - 2b, b - 30, -5)$.
\end{exercisesolution}

% Solution by Gemini 3.7 Flash (High)
\begin{exercisesolution}[Solution to \cref{exc:parametric_linear_system_solvability}]
\label{sol:parametric_linear_system_solvability}
\exinfo{Solution by Gemini 3.7 Flash (High). Computes the solvability hyperplane condition $b_1 - 4b_2 - b_3 = 0$ via row reduction.}
We form the augmented matrix and apply row operations:
\[
\begin{pmatrix}
0 & 7 & -2 & \big| & b_1 \\
-1 & 2 & -\frac{1}{2} & \big| & b_2 \\
4 & -1 & 0 & \big| & b_3
\end{pmatrix}
\xrightarrow{E_1 \leftrightarrow E_2}
\begin{pmatrix}
-1 & 2 & -\frac{1}{2} & \big| & b_2 \\
0 & 7 & -2 & \big| & b_1 \\
4 & -1 & 0 & \big| & b_3
\end{pmatrix}
\xrightarrow{-1 \cdot E_1 \to E_1}
\begin{pmatrix}
1 & -2 & \frac{1}{2} & \big| & -b_2 \\
0 & 7 & -2 & \big| & b_1 \\
4 & -1 & 0 & \big| & b_3
\end{pmatrix}
\]
Eliminate the first column of row 3:
\[
\xrightarrow{-4 \cdot E_1 + E_3 \to E_3}
\begin{pmatrix}
1 & -2 & \frac{1}{2} & \big| & -b_2 \\
0 & 7 & -2 & \big| & b_1 \\
0 & 7 & -2 & \big| & b_3 + 4b_2
\end{pmatrix}
\xrightarrow{-1 \cdot E_2 + E_3 \to E_3}
\begin{pmatrix}
1 & -2 & \frac{1}{2} & \big| & -b_2 \\
0 & 7 & -2 & \big| & b_1 \\
0 & 0 & 0 & \big| & b_3 + 4b_2 - b_1
\end{pmatrix}
\]
\textbf{Consistency condition:} The system possesses at least one solution if and only if the third row does not produce a contradiction ($0 = c$ with $c \ne 0$). Hence, the necessary and sufficient condition is:
\[
b_1 - 4b_2 - b_3 = 0.
\]
\textbf{Solution set when consistent:} When $b_1 - 4b_2 - b_3 = 0$, $z$ is a free variable. Setting $z = 14t$ for $t \in \mathbb{R}$:
\[
7y - 2z = b_1 \implies 7y = b_1 + 28t \implies y = \frac{1}{7}b_1 + 4t.
\]
From the first row:
\[
x = -b_2 + 2y - \frac{1}{2}z = -b_2 + 2\left(\frac{1}{7}b_1 + 4t\right) - 7t = \frac{2}{7}b_1 - b_2 + t.
\]
Thus, the solution set $S \subseteq \mathbb{R}^3$ is:
\[
S = \left\{ \begin{pmatrix} \frac{2}{7}b_1 - b_2 \\ \frac{1}{7}b_1 \\ 0 \end{pmatrix} + t \begin{pmatrix} 1 \\ 4 \\ 14 \end{pmatrix} \;\middle|\; t \in \mathbb{R} \right\}.
\]
\end{exercisesolution}

% Solution by Gemini 3.7 Flash (High)
\begin{exercisesolution}[Solution to \cref{exc:inconsistent_2x2_system}]
\label{sol:inconsistent_2x2_system}
\exinfo{Hybrid Solution (Gemini 3.7 Flash / Official Problem Sheet 4, Exercise 1). Proves the contrapositive geometric condition for $2 \times 2$ inconsistent systems using both determinant invertibility and echelon reduction.}
Let $A = \begin{pmatrix} a_{11} & a_{12} \\ a_{21} & a_{22} \end{pmatrix} \in \M_{2 \times 2}(\mathbb{R})$ and $b = \begin{pmatrix} b_1 \\ b_2 \end{pmatrix} \in \mathbb{R}^2$. 

Suppose for contradiction that neither row of $A$ is a scalar multiple of the other. This means that the row vectors $(a_{11}, a_{12})$ and $(a_{21}, a_{22})$ are linearly independent in $\mathbb{R}^2$. 

For a $2 \times 2$ matrix, linear independence of the rows is equivalent to
\[
\det(A) = a_{11}a_{22} - a_{12}a_{21} \ne 0.
\]
When $\det(A) \ne 0$, the matrix $A$ is invertible with inverse
\[
A^{-1} = \frac{1}{\det(A)} \begin{pmatrix} a_{22} & -a_{12} \\ -a_{21} & a_{11} \end{pmatrix}.
\]
Therefore, the system $Ax = b$ admits the unique solution $x = A^{-1}b$. But this contradicts the hypothesis that the system $Ax = b$ has no solutions.

Hence, the rows of $A$ must be linearly dependent, which means that one row of $A$ is a scalar multiple of the other.

\begin{ainote}
\textbf{Alternative Elementary Proof via Row Operations:}
We can also prove the contrapositive directly using only Gaussian elimination: suppose no row of $A = \begin{pmatrix} a & b \\ c & d \end{pmatrix}$ is a scalar multiple of the other.
\begin{itemize}
    \item If $a \ne 0$, subtract $\frac{c}{a} R_1$ from $R_2$ to obtain the second row $\bigl(0, \; d - \frac{bc}{a} \;\big|\; b_2 - \frac{c}{a}b_1\bigr)$. Since the rows are not proportional, $ad - bc \ne 0$, so the pivot $d - \frac{bc}{a} \ne 0$, which allows solving for both unknowns uniquely for any $(b_1, b_2)$.
    \item If $a = 0$, then $c \ne 0$ (otherwise row 1 is $0$ times row 2). Swapping rows reduces to the previous case.
\end{itemize}
Thus, if no row is a multiple of the other, the system is consistent for every right-hand side $b \in \mathbb{R}^2$.
\end{ainote}
\end{exercisesolution}

